\chapter*{Appendix B --- Phase 2 Plan}
\label{appendix:plan}
\addcontentsline{toc}{chapter}{Appendix B --- Phase 2 Plan}

This appendix details the Phase~2 plan in full, complementing the
narrative roadmap of Section~\ref{sec:future:roadmap}. Together,
the chapter and the appendix address Câu~5 of the grading
rubric~\cite{rubric} on all four dimensions: content, time,
personnel, and expected outcomes.

\section*{B.1 Deliverables per phase}

Each phase produces one or more deliverables that the next
phase depends on. The acceptance criterion for each phase is
the boolean condition under which that phase is considered
complete and the next phase can begin.

\begin{table}[H]
\centering
\caption{Phase 2 deliverables, expected outcomes, and acceptance
criteria.}
\label{tab:appendix:deliverables}
\footnotesize
\renewcommand{\arraystretch}{1.2}
\begin{tabularx}{\linewidth}{p{0.7cm}p{3.6cm}X}
\hline
\textbf{Ph.} & \textbf{Deliverable} & \textbf{Acceptance criterion} \\
\hline
P1 & Warehouse Layers 2--4 build
   & \texttt{python -m src.warehouse.builder} runs in
     $<$~30~s on local sample and emits ten \texttt{dim\_*}
     tables, three \texttt{fact\_*} views, and eight
     \texttt{agg\_*} tables; \texttt{notebook 05} asserts row
     counts match facts. \\
P2 & End-to-end Card~2 (Demand)
   & All five charts render; linked brushing works across all
     five; A3 entity panel opens on any route click; cold-load
     $<$~3~s; warm render $<$~50~ms. \\
P3 & Spatial card + pydeck
   & H3 hex layer renders at three cell sizes; Stop Demand
     Explorer renders 4{,}710 stops without UI lag. \\
P4 & School flagship card
   & Schools $\rightarrow$ Routes treemap renders;
     school-routeness classifier renders for all 163 routes;
     equity map renders at district resolution. \\
P5 & Financial + Network Health cards
   & P\&L quadrant clickable; QoS-1/2/3 scores computed on the
     nine routes available in the local GPS sample;
     threshold-$T$ slider snaps between the five pre-computed
     levels. \\
P6 & Routes card with Marey
   & Discovery and Detail modes both render; Marey diagram is
     interactive (zoom, hover trip); GroupSet matrix shows the
     four archetypes correctly. \\
P7 & Hub + entity-panel polish
   & Hub HTML renders inside Streamlit with no layout
     breakage; the A3 panel works for all six entity types;
     URL state survives a browser refresh. \\
P8 & Stories scrollytelling
   & All four stories scroll-driven; each scene deep-links into
     the matching card view; reading-progress bar visible. \\
P9 & Final polish + evaluation
   & School geocoding verified for $\geq$~50 of 64 schools;
     full 50-day GPS loaded if available; two-minute
     screen-recording demo produced; initial user evaluation
     with $\geq$~3 stakeholders per role completed. \\
\hline
\end{tabularx}
\end{table}

\section*{B.2 Risk register}

\begin{table}[H]
\centering
\caption{Top four risks for Phase 2.}
\label{tab:appendix:risk}
\footnotesize
\renewcommand{\arraystretch}{1.2}
\begin{tabularx}{\linewidth}{p{3.6cm}p{1.6cm}p{1.6cm}X}
\hline
\textbf{Risk} & \textbf{Likelihood} & \textbf{Impact} & \textbf{Mitigation} \\
\hline
Streamlit rerun latency exceeds 600~ms target
    & Medium & Medium
    & Pre-compute highlight sets server-side; only recompute
      the chart's colour column on rerun, not the full data. \\
Full 50-day GPS pull from Kaggle fails or delays
    & Medium & High
    & Demo with 1-day sample; document the seven charts that
      stay at SHOULD/WON'T tier until the pull completes. \\
School geocoding under 50/64 success rate
    & High & Medium
    & Ship with district-level approximation and the
      ``verified~\checkmark'' badge; do not block on full
      geocoding. \\
UMAP, parallel coordinates, or treemap fail to render
at 163 routes
    & Low & Low
    & Sanity-tested in EDA notebooks; all three render
      cleanly in Plotly at less than 500 routes. \\
\hline
\end{tabularx}
\end{table}

\section*{B.3 Personnel and time tracking}

Phase~2 is single-author. Each phase commit closes when the
phase's acceptance criterion in
Table~\ref{tab:appendix:deliverables} passes. The thesis
advisor reviews at the end of P2 (first end-to-end card) and
at the end of P6 (all cards rendering) before sign-off.

Time is tracked per phase, not per task, because every phase
is small enough for one developer to hold the full context.
The Gantt diagram in Section~\ref{sec:future:roadmap}
(Figure~\ref{fig:future:gantt}) shows the planned scheduling.

\section*{B.4 Post-Phase-2 user-validation study}

A formal user study is scheduled to begin one week after P9
completion. The study runs in two waves:

\begin{itemize}
\item \emph{Wave 1 (think-aloud sessions):} three or four
    members of each of the four user roles
    (Section~\ref{sec:proposed:problem}) work through a
    scripted set of analytical tasks while the moderator
    observes. Each session is approximately ninety minutes.
\item \emph{Wave 2 (free exploration):} the same participants
    are given the dashboard with no script for two weeks, with
    weekly check-ins. The outcome measure is the number of
    analytical findings the participant produces that the
    dashboard surfaced but the participant had not previously
    articulated.
\end{itemize}

The study is the basis for the ``Đánh giá và cải tiến''
Phase~2 deliverable required by the capstone topic
description~\cite{capstoneDescription}.
